\section{Related Work}
\subsection{Machine Learning Models as Classifiers in Educational Context}
The use of machine learning models as classifiers in educational contexts has been explored in various studies \cite{Jimenez-Macias2025, Zhan2025, Tong2025, Alhakbani2022}. From these studies, two main conclusions can be drawn.

First, certain models often outperform others in terms of accuracy for general classification tasks. Most prominently, ensemble methods---particularly Random Forest (RF)---consistently rank as top performers: Delgado et al. (2014) established RF as a "gold standard" across 179 classifiers \cite{Delgado2014}, a finding echoed in educational contexts by Jimenez-Macias et al. (2025) \cite{Jimenez-Macias2025}. Support Vector Machines and Artificial Neural Networks also prove effective for capturing non-linear relationships in educational data, with reported accuracies reaching 97.88\% and 99\%, respectively \cite{Zhan2025, Alhakbani2022}. Gradient Boosting variants such as XGBoost have also shown similar strength in identifying at-risk students \cite{Alhakbani2022, Delgado2014}. More recently, stacking ensemble models have been explored; these methods train heterogeneous base learners (e.g., KNN, Naive Bayes, RF, and XGBoost) and aggregate their outputs using a meta-learner such as Logistic Regression to mitigate individual model biases. For instance, a stacking framework achieved an accuracy of 98.53\% in predicting learning achievements on the XuetangX MOOC dataset \cite{Tong2025}.

Second, a critical consensus across current literature is the superiority of behavioral engagement data over demographic information for classification tasks. Online behavioral logs, including video watch counts, forum posts, and quiz participation, are significantly more influential in predicting academic risk than static factors like age, gender, or educational background. Jimenez-Macias et al. (2025) further refined this by demonstrating that time spent on tasks is a more sensitive predictor of performance than the raw number of attempts \cite{Jimenez-Macias2025}. Moreover, the distribution of grades across interactions significantly impacts model accuracy, as uniform grade distributions enhance the algorithm's ability to generalize patterns and make accurate predictions \cite{Jimenez-Macias2025}.

Methodological integrity in this context requires addressing class imbalance and interpretability. Educational datasets are frequently imbalanced, often containing significantly fewer failing students than passing ones \cite{Bujang2021, Tong2025}. To address this, researchers employed data-level solutions like SMOTE (Synthetic Minority Oversampling Technique) or Sequential Conditional GANs to balance class distributions and prevent model bias toward the majority class \cite{Bujang2021}. Furthermore, as models grow in complexity, tools like SHAP (SHapley Additive exPlanations) are increasingly used to provide "white-box" transparency, identifying which specific student behaviors drive the final prediction \cite{Tong2025}.

Despite these advancements, inconsistent preprocessing and feature scaling remain primary challenges, as they are often cited as the reason for varied performance reports for identical models across different studies \cite{Alhakbani2022}. Finally, the emerging paradigm of Tabular In-Context Learning (TabICL) using Large Language Models (LLMs) offers a novel interface for classification that can complement traditional frameworks like XGBoost and CatBoost, which have served as state-of-the-art benchmarks for tabular data for decades \cite{Wen2025}.

\subsection{Data aggregation}
It stands to reason that gathering high-quality and representative data is often a difficult endeavor \cite{Alias2024}. Alias et al. (2024) delineated a framework in which a list of hypothetically relevant features is enumerated and filtered via information gain. More specifically, the features taken into consideration include those relating to students' study hours, participation in extracurricular activities, hours of sleep, and mock exams practiced \cite{Suleiman2024}. Surveys for these features are deployed online as a questionnaire using the platform Google Forms.


\subsection{Performance metrics}
Common metrics to assess classification algorithms include accuracy, precision, recall, F1-score \cite{Alias2024, Zaky2023}, and the area under the receiver operating characteristic curve (AUC-ROC) \cite{Brzezinski2017, Yang2017}. These metrics provide insights into a model's ability to correctly classify instances, handle imbalanced datasets, and make informed decisions based on predictions \cite{Brzezinski2017}. Techniques for handling imbalanced datasets are crucial when working with decision tree (DT)-based models to avoid bias toward the majority class. Here, we focus on the Synthetic Minority Over-sampling Technique (SMOTE), specifically its variant ADASYN, which utilizes a weighted distribution to determine the learning difficulty of minority-class samples and generates synthetic data accordingly \cite{Chawla2011, Elreedy2019}. This approach helps to balance the dataset and improve the performance of decision tree models in predicting the minority class.